\documentclass[12pt,a4paper]{article}

% ── Paquetes ───────────────────────────────────────────────────────────────────
\usepackage[utf8]{inputenc}
\usepackage[spanish]{babel}
\usepackage{amsmath}
\usepackage{graphicx}
\usepackage{listings}
\usepackage{xcolor}
\usepackage{hyperref}
\usepackage{geometry}
\usepackage{tikz}
\usepackage{enumitem}
\usepackage{fancyhdr}

\geometry{margin=2.5cm}

% ── Estilo de código ──────────────────────────────────────────────────────────
\lstset{
    language=Python,
    basicstyle=\ttfamily\small,
    keywordstyle=\color{blue},
    stringstyle=\color{red},
    commentstyle=\color{green!50!black},
    numbers=left,
    numberstyle=\tiny,
    frame=single,
    breaklines=true,
    backgroundcolor=\color{gray!5},
}

% ── Configuración de página ───────────────────────────────────────────────────
\pagestyle{fancy}
\fancyhf{}
\fancyhead[L]{\textbf{MeteoApp} -- Guía de Exposición}
\fancyhead[R]{\thepage}
\setlength{\headheight}{14pt}

% ── Título ─────────────────────────────────────────────────────────────────────
\title{\Huge \textbf{MeteoApp}\\[0.3cm]
       \Large Dashboard Meteorológico Personal\\[0.3cm]
       \large Guía Técnica para Exposición}
\author{Valentina Rodriguez Sepulveda\\[0.2cm]
        \small ID: 1125789977\\[0.2cm]
        \small Introducción a la Programación -- Primer Semestre}
\date{\today}

\begin{document}

\maketitle
\thispagestyle{empty}
\newpage

\tableofcontents
\newpage

% ===============================================================================
\section{¿Qué es MeteoApp?}
% ===============================================================================

MeteoApp es una aplicación de escritorio desarrollada en Python que funciona como una \textbf{estación meteorológica digital}. Se conecta a internet para consultar datos climáticos reales de cualquier ciudad del mundo y los presenta al usuario a través de una interfaz gráfica amigable.

La aplicación utiliza la \textbf{API pública gratuita de Open-Meteo} para obtener datos como temperatura, humedad, velocidad del viento y precipitación. Toda la interacción se maneja a través de ventanas, botones y gráficas, haciendo que el usuario no necesite conocimientos técnicos para operarla.

\subsection{¿Qué puede hacer el usuario?}

\begin{itemize}[itemsep=4pt]
    \item \textbf{Iniciar sesión o registrarse} para tener su propio perfil.
    \item \textbf{Consultar el clima actual} de cualquier ciudad del mundo.
    \item \textbf{Guardar ciudades favoritas} para acceso rápido.
    \item \textbf{Ver gráficas de historial meteorológico} (temperatura y precipitación).
    \item \textbf{Configurar alertas de temperatura} que se activan cuando se superan umbrales definidos.
    \item \textbf{Cambiar entre unidades} °C/°F y km/h/mph \textbf{sin} volver a consultar la API.
    \item \textbf{Usar la aplicación sin internet}, gracias a un sistema de caché local.
\end{itemize}

% ===============================================================================
\section{Arquitectura del proyecto}
% ===============================================================================

\subsection{Estructura de carpetas}

El proyecto está organizado en una arquitectura de tres capas: \textbf{controlador (main)}, \textbf{vistas (views)} y \textbf{lógica (utils)}.

\begin{lstlisting}[frame=none, numbers=none, basicstyle=\ttfamily\footnotesize]
interfaz_grafica_flet/
|-- main.py
|-- requirements.txt
|-- utils/
|   |-- api_client.py
|   |-- data_manager.py
|   |-- settings.py
|   \-- chart_generator.py
|-- views/
|   |-- login_view.py
|   |-- home_view.py
|   |-- history_view.py
|   |-- cities_view.py
|   \-- alerts_view.py
|-- data/
|   |-- usuarios.csv
|   |-- ciudades.csv
|   |-- historial_clima.csv
|   \-- cache_clima.csv
\-- assets/
    \-- charts/
\end{lstlisting}

\noindent\textbf{Archivos y carpetas clave (breve explicación)}
\begin{itemize}
    \item \textbf{main.py}: punto de entrada y controlador de navegación entre pantallas.
    \item \textbf{utils/api\_client.py}: consulta Open-Meteo y traduce las claves JSON a español.
    \item \textbf{utils/data\_manager.py}: lectura y escritura de los CSV en \texttt{data/}.
    \item \textbf{utils/settings.py}: preferencias de unidades (°C/°F, km/h/mph).
    \item \textbf{utils/chart\_generator.py}: genera las gráficas PNG para \texttt{assets/charts/}.
    \item \textbf{views/}: contiene las pantallas de la app; \texttt{home\_view.py} es la principal.
    \item \textbf{data/}: CSV de persistencia (usuarios, ciudades, historial, cache).
\end{itemize}

% Nota para la exposición: una lista clara y sin comentarios en línea evita saltos de línea innecesarios al renderizar.

\subsection{Sugerencia: cómo explicar y demostrar en clase}
\begin{itemize}
    \item Empieza mostrando la interfaz principal (tarjeta de clima) y explica qué datos muestra.
    \item Demo rápida (pasos): iniciar sesión (usuario demo), buscar "Bogot\'a", seleccionar la ciudad, agregar a favoritas, cambiar unidad a °F, abrir Historial (7 días).
    \item Muestra el modo offline: desconectar internet y abrir la ciudad cacheada para que la audiencia vea el banner naranja.
    \item Explica brevemente la seguridad: las contrase\~nas se guardan como hash SHA‑256 (menciona que para producción se usar\'ian salt/bcrypt).
    \item Finaliza con arquitectura: separaci\'on clara entre vistas (UI) y \texttt{utils/} (l\'ogica y datos).
\end{itemize}

\subsection{Cómo ejecutar la aplicación en la exposición}

\noindent\textbf{Para la presentación:} se recomienda usar el ejecutable compilado para arranque instantáneo con doble clic, mientras se mantiene el código fuente abierto en el editor para explicar la arquitectura en tiempo real.

\begin{enumerate}[itemsep=4pt]
    \item \textbf{Primera vez (solo una vez):} hacer doble clic en \texttt{build\_windows.bat}. Este script crea un entorno virtual, instala todas las dependencias y compila \texttt{MeteoApp.exe} automáticamente.
    \item \textbf{Siempre:} doble clic en \texttt{MeteoApp.exe} para abrir la aplicación.
    \item \textbf{En paralelo (durante la exposición):} abrir el código fuente (\texttt{main.py}, \texttt{utils/}, \texttt{views/}) en VS Code o cualquier editor. El ejecutable lee y escribe los archivos \texttt{data/*.csv} en tiempo real, por lo que el público puede ver cómo se crean y actualizan los registros CSV en el editor mientras se usa la aplicación.
\end{enumerate}

\noindent\textbf{Nota técnica:} \texttt{utils/chart\_generator.py} usa el backend \texttt{Agg} de matplotlib (sin interfaz gráfica) para generar los PNG. Esto funciona correctamente tanto en el ejecutable como al correr desde el código fuente con \texttt{python main.py}.

\newpage
\subsection{Diagrama de flujo de datos}

\begin{figure}[h]
    \centering
    \begin{tikzpicture}[node distance=2cm, auto]
        % Nodos
        \node (user) [draw, rounded corners, fill=blue!20, minimum width=3cm, minimum height=1cm] {Usuario};
        \node (view) [draw, rounded corners, fill=green!20, minimum width=3cm, minimum height=1cm, below of=user] {Vista (view/*.py)};
        \node (api)  [draw, rounded corners, fill=orange!20, minimum width=3cm, minimum height=1cm, below of=view] {API Client (utils/api\_client.py)};
        \node (openmeteo) [draw, rounded corners, fill=red!20, minimum width=3cm, minimum height=1cm, below of=api] {Open-Meteo API};
        \node (csv)  [draw, rounded corners, fill=purple!20, minimum width=3cm, minimum height=1cm, right of=api, xshift=4cm] {CSV Files (data/*.csv)};

        % Flechas
        \draw[->, thick] (user) -- (view) node[midway, right] {clic};
        \draw[->, thick] (view) -- (api) node[midway, right] {consulta};
        \draw[->, thick] (api) -- (openmeteo) node[midway, right] {HTTP};
        \draw[->, thick] (openmeteo) -- (api) node[midway, left] {JSON};
        \draw[->, thick] (api) -- (view) node[midway, left] {datos};
        \draw[->, thick] (view) -- (csv) node[midway, above] {guardar};
        \draw[->, thick] (csv) -- (view) node[midway, below] {leer};
        \draw[->, thick] (view) -- (user) node[midway, left] {pantalla};
    \end{tikzpicture}
    \caption{Flujo de datos en MeteoApp. El usuario interactúa con la vista, que a su vez consulta la API o los archivos CSV locales.}
\end{figure}

\newpage
% ===============================================================================
\section{main.py -- El punto de entrada}
% ===============================================================================

El archivo \texttt{main.py} es el \textbf{corazón de la navegación}. Cuando se ejecuta la aplicación (ya sea con \texttt{python main.py} o mediante el ejecutable \texttt{MeteoApp.exe}), ocurre lo siguiente:

\begin{enumerate}[itemsep=4pt]
    \item Se crea una ventana de 960x680 píxeles con tema azul y fuente Roboto.
    \item Se llama a \texttt{mostrar\_inicio\_sesion()}, que carga la primera pantalla: el inicio de sesión.
    \item Cada función de navegación (\texttt{mostrar\_inicio}, \texttt{mostrar\_historial}, \texttt{mostrar\_ciudades}, \texttt{mostrar\_alertas}) sigue el mismo patrón:
    \begin{enumerate}
        \item Limpiar la pantalla anterior.
        \item Crear una instancia de la vista correspondiente.
        \item Pasarle \"callbacks\" (funciones que se ejecutarán cuando el usuario haga clic en un botón).
        \item Añadir la vista a la página.
    \end{enumerate}
\end{enumerate}

\textbf{Concepto clave: callbacks.} Una vista no sabe cómo navegar a otra pantalla; recibe una función (callback) que main.py le entrega. Cuando el usuario hace clic en "Historial", la vista ejecuta el callback y main.py se encarga del resto.

% ===============================================================================
\section{La carpeta utils/ -- La lógica interna}
% ===============================================================================

La carpeta \texttt{utils/} contiene todo el código que \textbf{no dibuja nada en pantalla}. Es la parte \"invisible\" pero fundamental: procesa datos, habla con la API, lee archivos y genera gráficas.

\subsection{utils/api\_client.py -- Comunicación con la API}

Este archivo es el \textbf{mensajero} de la aplicación. Se encarga de enviar peticiones a la API de Open-Meteo y traducir las respuestas.

\subsubsection*{¿Qué es una API?}

API son las siglas de \textit{Application Programming Interface} (Interfaz de Programación de Aplicaciones). En términos sencillos:

\begin{quote}
Una API es como un \textbf{mesero en un restaurante}: tú le pides algo (\"quiero el clima de Bogotá\"), él va a la cocina (el servidor de Open-Meteo), trae tu pedido y te lo entrega en un formato que puedes entender (JSON). Tú no necesitas saber cómo se cocinó; solo recibes lo que pediste.
\end{quote}

\subsubsection*{Los tres endpoints de Open-Meteo}

\begin{itemize}
    \item \textbf{Geocodificación} -- Convierte el nombre de una ciudad en coordenadas (latitud y longitud). Se usa cuando el usuario escribe "Bogotá" y la app necesita saber dónde está.
    
    \item \textbf{Pronóstico} -- Devuelve el clima actual y el pronóstico del día. Se usa en la pantalla principal para mostrar temperatura, humedad, viento, etc.
    
    \item \textbf{Archivo histórico} -- Devuelve datos meteorológicos de fechas pasadas. Se usa en el historial para generar las gráficas.
\end{itemize}

\subsubsection*{Traducción de claves JSON}

La API de Open-Meteo responde en inglés (por ejemplo, \texttt{temperature\_2m}). api\_client.py \textbf{traduce} estas claves a español antes de entregarlas a las vistas:

\begin{lstlisting}[frame=none, numbers=none, basicstyle=\ttfamily\small]
"temperature_2m"           --> "temperatura"
"apparent_temperature"     --> "sensacion_termica"
"relative_humidity_2m"    --> "humedad"
"wind_speed_10m"          --> "viento"
"weathercode"             --> "codigo_clima" + "descripcion" + "icono"
"temperature_2m_max"      --> "temp_max"
"temperature_2m_min"      --> "temp_min"
"precipitation_sum"       --> "precipitacion"
\end{lstlisting}

Así, las vistas nunca ven los nombres en inglés de la API; trabajan con nombres en español que son más fáciles de entender.

\subsubsection*{¿Por qué 2 metros y 10 metros?}

Los estándares meteorológicos mundiales establecen que:
\begin{itemize}
    \item La \textbf{temperatura} se mide a 2 metros del suelo (altura de un termómetro estándar).
    \item El \textbf{viento} se mide a 10 metros del suelo (altura de un anemómetro en una estación meteorológica).
\end{itemize}

\subsubsection*{Manejo de errores de red}

Cuando la aplicación no puede comunicarse con la API, el código distingue entre tres tipos de error:

\begin{itemize}
    \item \texttt{requests.Timeout} -- El servidor no respondió en 10 segundos.
    \item \texttt{requests.ConnectionError} -- No hay conexión a internet.
    \item \texttt{requests.HTTPError} -- El servidor respondió con un código de error (ej. 404).
\end{itemize}

En todos los casos, la función devuelve \texttt{None} (sin datos), permitiendo que la vista decida si mostrar caché o un mensaje de error.

\newpage
\subsection{utils/data\_manager.py -- Persistencia de datos}

Este archivo maneja la \textbf{lectura y escritura de archivos CSV} usando la librería pandas.

\subsubsection*{¿Qué es persistencia de datos?}

Persistencia significa que los datos \textbf{no se pierden cuando se cierra la aplicación}. En MeteoApp, los datos se guardan en archivos CSV (Comma-Separated Values), que son tablas de texto plano donde cada fila es un registro y cada columna es un campo.

\begin{lstlisting}[frame=none, numbers=none, basicstyle=\ttfamily\small]
# Ejemplo de usuarios.csv (visto como tabla)
| id | username | password_hash                   | fecha_registro |
|----|----------|---------------------------------|----------------|
| 1  | Valen    | a1b2c3d4e5f6... (64 caracteres) | 2026-05-30     |
| 2  | Juan     | 9f8e7d6c5b4a... (64 caracteres) | 2026-06-01     |
\end{lstlisting}

\subsubsection*{Los cuatro archivos CSV}

\begin{enumerate}
    \item \textbf{usuarios.csv} -- Almacena las cuentas de usuario: ID, nombre de usuario, hash de contraseña y fecha de registro.
    \item \textbf{ciudades.csv} -- Almacena las ciudades favoritas de cada usuario, sus coordenadas, umbrales de alerta (en °C y °F) y el último clima consultado (para modo offline).
    \item \textbf{historial\_clima.csv} -- Almacena datos históricos descargados de la API para evitar consultas repetidas. Se purga automáticamente después de 365 días.
    \item \textbf{cache\_clima.csv} -- Almacena el último clima consultado de cada ciudad (incluso las no favoritas), para que la aplicación funcione sin internet.
\end{enumerate}

\subsection{utils/settings.py -- Preferencias de unidades}

Este archivo mantiene en memoria las preferencias del usuario durante la sesión:

\begin{itemize}
    \item \textbf{Temperatura:} °C o °F.
    \item \textbf{Velocidad:} km/h o mph.
\end{itemize}

Las variables son \textbf{globales} (existen fuera de cualquier clase) porque la preferencia debe mantenerse aunque el usuario navegue entre ventanas. Si cambia °C → °F en la pantalla principal, al ir al historial debe seguir en °F.

\textbf{Importante:} Las conversiones se hacen en tiempo real sin volver a consultar la API. La fórmula es:
\begin{align*}
    F &= C \times \frac{9}{5} + 32\\
    \text{mph} &= \text{km/h} \times 0.621371
\end{align*}

\subsection{utils/chart\_generator.py -- Generación de gráficas}

Este archivo usa \texttt{matplotlib} para crear gráficas profesionales en formato PNG:

\begin{itemize}
    \item \textbf{chart\_temperatura()} -- Genera una gráfica de líneas con la temperatura máxima (rojo) y mínima (azul) diaria, con relleno entre ambas.
    \item \textbf{chart\_precipitacion()} -- Genera una gráfica de barras con la lluvia acumulada por día en milímetros (1 mm = 1 litro de agua por metro cuadrado).
\end{itemize}

Las imágenes se guardan en \texttt{assets/charts/} y se convierten a base64 para mostrarlas en la interfaz de Flet.

% ===============================================================================
\section{Las vistas (views/) -- La interfaz gráfica}
% ===============================================================================

Cada archivo en \texttt{views/} representa una \textbf{pantalla} de la aplicación. Todas siguen la misma estructura: una clase con un método \texttt{build()} que construye y devuelve la interfaz.

\subsection{views/login\_view.py -- Inicio de sesión}

\begin{itemize}
    \item Muestra un formulario con campos de usuario y contraseña.
    \item Permite alternar entre modo "Iniciar sesión" y "Registrarse".
    \item Valida que los campos no estén vacíos y que la contraseña tenga al menos 4 caracteres.
    \item Las contraseñas se almacenan como \textbf{hash SHA-256}, nunca en texto plano.
\end{itemize}

\subsection{views/home\_view.py -- Pantalla principal}

Es el archivo \textbf{más grande y complejo} (674 líneas). Contiene:

\begin{itemize}
    \item \textbf{Buscador de ciudades} con sugerencias (hasta 5 resultados) y selección mediante diálogo.
    \item \textbf{Tarjeta del clima} que muestra temperatura, sensación térmica, humedad, viento e ícono del clima.
    \item \textbf{Botones de unidad} para alternar entre °C/°F y km/h/mph.
    \item \textbf{Banner de alerta} que se activa cuando la temperatura supera los umbrales configurados.
    \item \textbf{Modo offline} que muestra el último clima cacheado con un aviso naranja.
    \item \textbf{Diálogo de favoritos} que pregunta si se desea agregar la ciudad consultada a la lista de favoritas.
    \item \textbf{Íconos del clima} generados \textbf{con Pillow} en tiempo real: sol, nubes, lluvia, nieve, tormenta, niebla.
\end{itemize}

\subsection{views/history\_view.py -- Historial meteorológico}

\begin{itemize}
    \item Permite seleccionar una ciudad favorita, una fecha de inicio y un número de días (7 a 183).
    \item Consulta primero la \textbf{caché local} (historial\_clima.csv); si no hay suficientes datos, llama a la API.
    \item Genera dos gráficas con matplotlib y las muestra en pantalla.
    \item Calcula estadísticas: promedio de temperatura máxima, promedio de mínima, promedio de lluvia y lluvia total del período.
    \item La consulta se ejecuta en un \textbf{hilo separado (threading)} para no congelar la interfaz mientras se descargan datos.
\end{itemize}

\subsection{views/cities\_view.py -- Ciudades favoritas}

\begin{itemize}
    \item Muestra la lista de ciudades favoritas del usuario con checkboxes.
    \item Permite agregar nuevas ciudades (las geocodifica primero para obtener coordenadas).
    \item Permite eliminar ciudades seleccionadas.
    \item Muestra el conteo de ciudades guardadas.
\end{itemize}

\subsection{views/alerts\_view.py -- Alertas de temperatura}

\begin{itemize}
    \item Permite seleccionar una ciudad favorita y definir umbrales de temperatura máxima y mínima.
    \item Los valores se ajustan con sliders que se validan entre sí (máximo no puede ser menor que mínimo).
    \item Cada alerta se guarda en °C y también se precalcula en °F para mostrarla rápido si el usuario cambia de unidad.
    \item Muestra un resumen de todas las alertas activas.
\end{itemize}

\newpage
% ===============================================================================
\section{Seguridad: Hash SHA-256 de contraseñas}
% ===============================================================================

MeteoApp \textbf{nunca} guarda las contraseñas en texto plano. En su lugar, las convierte en un \textbf{hash} de 64 caracteres hexadecimales utilizando el algoritmo SHA-256.

\subsection{¿Qué es un hash?}

Un hash es una \textbf{huella digital}: una función matemática que toma un texto de cualquier tamaño y produce un valor fijo (64 caracteres). Las propiedades clave son:

\begin{enumerate}
    \item \textbf{Irreversible:} No se puede obtener la contraseña original a partir del hash.
    \item \textbf{Determinístico:} La misma contraseña siempre produce el mismo hash.
    \item \textbf{Unidireccional:} Cambiar un solo carácter de la contraseña produce un hash completamente diferente.
\end{enumerate}

\subsection{¿Cómo funciona en la aplicación?}

\begin{lstlisting}[language=Python, basicstyle=\ttfamily\small]
import hashlib

def _codificar_contrasena(password: str) -> str:
    """Convierte la contrasena en una huella digital SHA-256."""
    return hashlib.sha256(password.encode()).hexdigest()
\end{lstlisting}

El proceso es:
\begin{enumerate}
    \item El usuario escribe su contraseña (ej. "MiClave2026").
    \item La aplicación calcula su hash SHA-256 (64 caracteres hexadecimales).
    \item Guarda \textbf{solo el hash} en usuarios.csv.
    \item Al iniciar sesión, vuelve a calcular el hash de la contraseña ingresada y lo compara con el guardado.
\end{enumerate}

\begin{quote}
\textbf{Ejemplo visual:}
``hola'' → SHA-256 → ``1406e...b0d3'' (64 caracteres)\\
``Hola'' → SHA-256 → ``b221d...a4f1'' (completamente diferente por la mayúscula)
\end{quote}

% ===============================================================================
\section{Librerías utilizadas}
% ===============================================================================

\begin{tabular}{|l|l|l|}
\hline
\textbf{Librería} & \textbf{Versión} & \textbf{¿Para qué se usa?} \\
\hline
Flet & 0.28.3 & Interfaz gráfica (botones, ventanas, menús) \\
Requests & -- & Consultas HTTP a la API de Open-Meteo \\
Pandas & -- & Manejo de datos tabulares (CSV, DataFrames) \\
Matplotlib & -- & Generación de gráficas científicas \\
Pillow & -- & Creación de íconos del clima en tiempo real \\
\hline
\end{tabular}

\subsection*{¿Cómo se instalan?}

\begin{lstlisting}[frame=none, numbers=none, basicstyle=\ttfamily\small]
pip install -r requirements.txt
\end{lstlisting}

El archivo \texttt{requirements.txt} contiene exactamente:
\begin{lstlisting}[frame=none, numbers=none, basicstyle=\ttfamily\small]
flet==0.28.3
pandas
requests
matplotlib
Pillow
\end{lstlisting}

% ===============================================================================
\section{Resumen técnico rápido}
% ===============================================================================

\begin{itemize}[itemsep=4pt]
    \item \textbf{Lenguaje:} Python 3.10+
    \item \textbf{Framework gráfico:} Flet 0.28.3 (basado en Flutter)
    \item \textbf{API externa:} Open-Meteo (gratuita, sin clave)
    \item \textbf{Persistencia:} 4 archivos CSV gestionados con pandas
    \item \textbf{Seguridad:} Hash SHA-256 para contraseñas
    \item \textbf{Gráficas:} Matplotlib con backend sin GUI (Agg)
    \item \textbf{Íconos:} Generados dinámicamente con Pillow
    \item \textbf{Asincronía:} threading para consultas de historial sin congelar la UI
    \item \textbf{Arquitectura:} Separación en vistas (views/) y lógica (utils/)
    \item \textbf{Modo offline:} Caché local en CSV con detección automática
\end{itemize}

\end{document}
